%========================================================= §3 ATTACK STEPS
\section{Attack Steps \& Execution}\label{sec:steps}

This section answers requirement~(a): the concrete steps of the attack, with the tool's
own output as evidence and captioned placeholders for the demonstration screenshots. The
whole chain is driven by one command, for example:

\begin{logbox}
python attacker/main.py --target 192.168.0.116 --ports 21 22 80 81 \
       --adapter "\Device\NPF_{...}"
\end{logbox}

\subsection{Phase 1 — TCP SYN port scanning}
The scanner sends a raw SYN to each target port and classifies the reply: a SYN/ACK means
\textsc{open} (and no ACK is returned, keeping the connection half-open); a RST means
\textsc{closed}; silence past the timeout means \textsc{filtered}. Every open/closed reply
is also retained as OS-fingerprint evidence.

\begin{logbox}
[PHASE 1] TCP SYN Port Scanning  ->  192.168.0.116
   21  OPEN                           TTL=128 WIN=65535
   80  OPEN                           TTL=128 WIN=65535
   22  NO RESPONSE / POSSIBLY FILTERED
   81  NO RESPONSE / POSSIBLY FILTERED
\end{logbox}

\shotplaceholder{Attacker console — Phase 1 SYN scan output showing OPEN / CLOSED /
FILTERED port states with TTL and window values.}

\subsection{Phase 2 — service banner grabbing}
For each open port the grabber opens a full TCP connection and reads the service's
identifying banner; for HTTP it sends a \texttt{HEAD} probe and parses the \texttt{Server}
header. Banners feed both the report and the cross-layer fingerprint stage.

\begin{logbox}
[PHASE 2] Service Banner Grabbing
   21  FTP   -> 220 LabFTP Server Ready
   80  HTTP  -> Server: VulnHTTP/1.0 (Lab Demo)
\end{logbox}

\subsection{Phase 3 — OS fingerprinting}
The retained response features are scored by the probabilistic engine, which emits a most
likely OS family, a confidence percentage, a coarse version tier, and the evidence behind
the score. The engine and its seven strategies are detailed in Section~\ref{sec:fp}; a
representative verdict:

\begin{logbox}
[PHASE 3] OS Fingerprinting
   Likely OS  : Windows family        Confidence : 62%
   Version    : Windows 10 / 11 tier
   Evidence   : TTL=128 | WIN=65535 | IP-ID strict-incremental [27057,27059,27062]
              : TCP timestamp clock ~1000 Hz | ECN ignored | Windows option order
\end{logbox}

\subsection{Phase 4 — OS-aware payload selection}
Before any exploit runs, \texttt{os\_payloads.py} adapts the payloads to the fingerprinted
OS: a Windows verdict selects \texttt{windows\textbackslash win.ini} and \texttt{whoami};
a Linux verdict selects \texttt{/etc/passwd} and \texttt{id}. This closes the loop between
reconnaissance and weaponisation — the same binary attacks a Windows and a Linux host
differently, automatically.

\subsection{Phase 5 — exploit battery}
Each applicable exploit then executes. The FTP suite and HTTP traversal are pure-socket;
SSH uses \texttt{paramiko} for key exchange only, with our own attack logic on top.

\begin{itemize}
\item \textbf{FTP anonymous login} — authenticates as \texttt{anonymous}, then
\texttt{PWD}/\texttt{LIST} to prove data access.
\item \textbf{FTP write $\rightarrow$ web-shell $\rightarrow$ RCE} — uploads
\texttt{shell.py} into the writable \texttt{uploads/} directory, requests
\texttt{/uploads/shell.py?cmd=whoami} over HTTP to execute it, then deletes the artefact.
\item \textbf{FTP bounce} — attempts a \texttt{PORT} pivot; on a modern host this is
correctly rejected (see Section~\ref{sec:results}).
\item \textbf{HTTP directory traversal} — reads the OS-selected sensitive file across
several encodings (\texttt{..\textbackslash}, \texttt{../}, URL-encoded).
\item \textbf{SSH weak credentials} — brute-forces a small bounded wordlist, then runs
silent recon (quiet mode) or drops desktop trophy files and a spoken alert (loud mode).
\end{itemize}

\begin{logbox}
[EXPLOIT] FTP Web-Shell RCE (port 21+80)
  Status : SUCCESS   whoami -> homed01\turjo   (artefact cleaned up)
[EXPLOIT] HTTP Directory Traversal (port 80)
  Status : SUCCESS   read windows\win.ini  ([fonts] Arial=arial.ttf ...)
[EXPLOIT] FTP Bounce (port 21)
  Status : MITIGATED  [WinError 10053] connection aborted by host WFP
\end{logbox}

\shotplaceholder{Attacker console — Phase 5 exploit battery: FTP anonymous listing,
web-shell RCE output, and HTTP traversal file contents.}

\shotplaceholder{Victim desktop (SSH loud mode) — the injected \texttt{HACKED.py} script
and the \texttt{.bak}/text trophy file left on the Desktop with the attacker's message.
(The spoken voice alert is triggered on the victim's speakers and cannot be shown in a
still image.)}
